\section{A uniform spectral normal form}
\label{sec:uniform-normal-v96}
We compute the whole-model observation ball, not just the image of a
chosen root-splitting path. For a fixed centre in $\mathfrak I$ define
\[
 \langle H,H'\rangle_{I,\theta}
   =\sum_jw_j\sum_e{H_{je}H'_{je}\over P_{\theta,je}}.
\]
The column derivatives in
$\eta=(\dot\alpha,\dot u_1,\dot u_2,\dot v_1,\dot v_2)$ have entry
sum zero, so $\eta$ has five free scalar coordinates. Define
\begin{align}
 p_f(z)&=-a(z-r)^2-bz(z-r)-Xz,\nonumber\\
 p_g(z)&=-c(z-s)^2-Y(z-s),\label{eq:jets-v96}\\
 \dot K&=\alpha u_1v_1^{\mathsf T}p_f
              +\gamma u_2v_2^{\mathsf T}p_g\nonumber\\
 &\quad+f\{\dot\alpha u_1v_1^{\mathsf T}
        +\alpha\dot u_1v_1^{\mathsf T}+\alpha u_1\dot v_1^{\mathsf T}\}
       +g\{-\dot\alpha u_2v_2^{\mathsf T}
        +\gamma\dot u_2v_2^{\mathsf T}+\gamma u_2\dot v_2^{\mathsf T}\},\nonumber\\
 \dot q&=\alpha p_f+\gamma p_g+\dot\alpha(f-g),\nonumber\\
 S_\theta(a,b,c,\eta,X,Y)_j
   &={\dot K(T_j)-P_{\theta,j}\dot q(T_j)\over q(T_j)}.\label{eq:score-v96}
\end{align}
Here $a,X,Y\ge0$, whereas $b,c,\eta$ are free. The first inequality
is the endpoint condition; the other two describe real root splitting.
These formulas allow $u_1\ne u_2$.

Let $L_\theta$ be the seven-column matrix of the free directions
$(b,c,\eta)$, and let $D_\theta$ be the three-column matrix of the
directions $(a,X,Y)$. Let $W_\theta$ be the diagonal matrix with entries
$w_j/P_{\theta,je}$. Define
\begin{align}
 Q_\theta&=D_\theta^{\mathsf T}W_\theta D_\theta
 -D_\theta^{\mathsf T}W_\theta L_\theta
 (L_\theta^{\mathsf T}W_\theta L_\theta)^{-1}
 L_\theta^{\mathsf T}W_\theta D_\theta,\label{eq:Q-v96}\\
 J_\theta(X,Y)&=\min_{a\ge0}(a,X,Y)Q_\theta(a,X,Y)^{\mathsf T},\nonumber\\
 \kappa_x(\theta)&=\min_{Y\ge0}J_\theta(1,Y),\qquad
 \kappa_y(\theta)=\min_{X\ge0}J_\theta(X,1).\label{eq:kappas-v96}
\end{align}
The following lemma proves, in particular, that the inverse in
\eqref{eq:Q-v96} exists even at rank-one observations.

\begin{lemma}[Uniform feasible coercivity]\label{lem:coercive-v96}
The matrix $L_\theta$ has rank seven, $Q_{\theta,00}>0$, and
$\kappa_x,\kappa_y$ are positive and continuous on $\mathfrak I$.
For every compact $H\Subset\mathfrak I$ there is $c_H>0$ such that
\begin{equation}\label{eq:coercivity-v96}
 \|S_\theta(a,b,c,\eta,X,Y)\|_{I,\theta}
 \ge c_H\|(a,b,c,\eta,X,Y)\|
 \quad(a,X,Y\ge0,\ \theta\in H).
\end{equation}
In particular $J_\theta(X,Y)\ge c'_H(X+Y)^2$ on the nonnegative
quadrant. Every minimum in \eqref{eq:kappas-v96} is attained, with
minimizing jet variables uniformly bounded on $H$ after fixing $X=1$
or $Y=1$.
\end{lemma}
\begin{proof}
Every feasible jet is realized by admissible root and parameter arcs.
For $t>0$ choose a real triple $y$ with $\sum_i y_i=0$ and
$\frac12\sum_i y_i^2=Y$, and use roots
\begin{equation}\label{eq:arcs-v96}
 ta,\quad r+tb/2\pm\sqrt{tX},\quad
 s+tc/3+\sqrt t\,y_i\quad(1\le i\le3),
\end{equation}
together with parameter increment $t\eta$. Positivity and the interior
margins at $r,s$ make these genuine closed-model alternatives, for
small $t$, including when some roots coincide. Their observation
increment is $tS_\theta+O(t^{3/2})$.

If a feasible jet has zero score, Lemma~\ref{lem:uniform-inverse-v96}
applied to these arcs makes all polynomial and non-root derivatives
zero. Thus $p_f=p_g=0$ and $\eta=0$. Evaluation at $z=0$, followed by
coefficient comparison in \eqref{eq:jets-v96}, gives
$a=b=c=X=Y=0$. The score therefore has no nonzero zero on the closed
feasible cone. Compactness of its intersection with the unit sphere,
and then of $H$, proves \eqref{eq:coercivity-v96}.
This also proves that the free columns are independent. If the endpoint
column were in their span, a feasible jet with $a=1$ would have zero
score; hence $Q_{00}>0$.

Coercivity proves existence and boundedness of minimizing jets.
The upper bound is obtained by setting all other variables to zero;
the lower bound follows from \eqref{eq:coercivity-v96}. A convergent
sequence of centres has bounded minimizing jets, so a subsequence and
continuity of the score prove lower semicontinuity of either minimum.
Evaluation at a fixed minimizing jet proves upper semicontinuity.
This proves continuity and all remaining assertions.
\end{proof}

Let $\mathcal L_\theta$ be the compact set of triples of ordered root
clusters
\begin{equation}\label{eq:leading-set-v96}
 \bigl(\{0\},\{-\sqrt X,\sqrt X\},Z(z^3-Yz-Z)\bigr),\quad
 X,Y\ge0,\quad J_\theta(X,Y)\le4,\quad27Z^2\le4Y^3.
\end{equation}
Its metric is the maximum of the three ordered-cluster metrics. For
small $t$, define $\mathcal R_\theta(t)$ by taking every competitor in
the observation ball, grouping its roots near $0,r,s$, subtracting
these centres, and dividing by $\sqrt t$. Lemma~\ref{lem:uniform-inverse-v96}
ensures that these groups have sizes one, two and three, uniformly on
compact subsets of $\mathfrak I$.

\begin{theorem}[Uniform leading set and linear remainder]
\label{thm:uniform-normal-v96}
For every compact $H\Subset\mathfrak I$ there exist $C_H,t_H>0$ such
that, for $\theta\in H$ and $0<t<t_H$,
\begin{equation}\label{eq:uniform-Hausdorff-v96}
 d_H\bigl(\mathcal R_\theta(t),\mathcal L_\theta\bigr)
       \le C_H\sqrt t.
\end{equation}
Consequently the entire closed-model modulus satisfies
\begin{align}
 \omega_\theta(t)&=C(\theta)\sqrt t+O_H(t),\label{eq:uniform-law-v96}\\
 C(\theta)&=\max\left\{\sqrt2\,\kappa_x(\theta)^{-1/4},\,
              2\sqrt{2/3}\,\kappa_y(\theta)^{-1/4}\right\}.
                                                   \label{eq:C-v96}
\end{align}
The error is uniform across channel-rank changes within $H$.
All points of the leading set, including diameter-extremizing pairs,
are realized by actual stochastic alternatives. The function $C$ is
positive and continuous on $\mathfrak I$.
\end{theorem}
\begin{proof}
We give quantitative necessity and sufficiency. Positivity on $H$ and
\eqref{eq:uniform-inverse-v96} imply that every radius-$t$ competitor
has coefficient and non-root parameter increments $O_H(t)$.
Coprime factorization of its first component, at $0$ and $r$, gives
a simple root $ta$ and pair sum $2r+tb$, with $a\ge0$ and bounded
$a,b$. Its pair is exactly $r+tb/2\pm\sqrt{tX}$ with $X\ge0$ bounded.
The second component has mean $s+tc/3$ and roots
$s+tc/3+\sqrt t\,y_i$, where $\sum_i y_i=0$.
The first two symmetric coefficients give $\sum_i y_i^2=O_H(1)$.
Thus $Y=\frac12\sum_i y_i^2$ and $Z=\prod_i y_i$ are bounded and
satisfy $Y\ge0$, $27Z^2\le4Y^3$. The parameter increment is $t\eta$
with $\eta$ uniformly bounded. These are exact coordinates for the
competitor, not a choice of an approximating arc.

Expanding its component polynomials gives \eqref{eq:jets-v96} at
order $t$. The only possible order-$t^{3/2}$ term needed here is the
centred triple product; all remaining errors are bounded by that order.
The rational probability map has uniformly bounded derivatives on
these coefficient neighbourhoods. Hence, uniformly in all the
competitors under consideration,
\begin{equation}\label{eq:uniform-expansion-v96}
 F_T(\xi)=P_\theta+tS_\theta+O_H(t^{3/2}),\qquad
 h_w(F_T(\xi),P_\theta)
       ={t\over2}\|S_\theta\|_{I,\theta}+O_H(t^{3/2}).
\end{equation}
The second assertion follows by Taylor expansion of the vector of
square roots and the reverse triangle inequality for its Euclidean
norm; it does not require a nonzero score.

For necessity, ball membership gives $\|S_\theta\|_{I,\theta}
\le2+M_H\sqrt t$. Let $\rho=2/(2+M_H\sqrt t)$. Scaling all first-order
jet variables by $\rho$, and $Z$ by $\rho^{3/2}$, produces a point of
\eqref{eq:leading-set-v96}. Its root clusters differ from the
competitor's rescaled clusters by $O_H(\sqrt t)$: the centres contribute
$\sqrt t(a,b/2,c/3)$ and the radial change multiplies bounded centred
roots by $\sqrt\rho=1+O_H(\sqrt t)$. This proves one directed
Hausdorff estimate.

For sufficiency, take any point of \eqref{eq:leading-set-v96} and
choose minimizing nuisance variables and $a\ge0$. By
Lemma~\ref{lem:coercive-v96} these are uniformly bounded, because the
score is at most two. For a fixed sufficiently large $M'_H$, put
$\rho=1-M'_H\sqrt t>0$ and use the actual arcs \eqref{eq:arcs-v96}
with first-order variables multiplied by $\rho$, and centred triple
roots by $\sqrt\rho$. Their score is multiplied by $\rho$.
The uniform expansion \eqref{eq:uniform-expansion-v96} then gives
$h_w\le t(\rho+O_H(\sqrt t))\le t$. Their rescaled roots differ
from the prescribed point by $O_H(\sqrt t)$. All stochastic and root
constraints hold uniformly for small $t$. This proves the reverse
Hausdorff estimate, including boundary points of the Fisher ball,
without exchanging a pointwise limit with a family supremum.

It remains to compute the diameter. Homogeneity gives
$X_{\max}=2/\sqrt{\kappa_x}$ and
$Y_{\max}=2/\sqrt{\kappa_y}$. The double-cluster diameter is
$\sqrt{X_{\max}}$. A sorted zero-sum triple with
$\sum_i y_i^2\le2Y_{\max}$ satisfies
\[
 -2\sqrt{Y_{\max}/3}\le y_1\le0,\quad
 |y_2|\le\sqrt{Y_{\max}/3},\quad
 0\le y_3\le2\sqrt{Y_{\max}/3}.
\]
The outer-coordinate bounds follow from Cauchy--Schwarz and zero sum;
the middle bound follows by minimizing the norm subject to ordering.
The resulting diameter $2\sqrt{Y_{\max}/3}$ is attained by
$(-2h,h,h)$ and $(-h,-h,2h)$, $h=\sqrt{Y_{\max}/3}$.
They have the same $Y$ and opposite $Z$, hence the same leading score.
The maximum-cluster metric has diameter equal to its largest marginal
diameter; no independence of the cluster coordinates is assumed.
This proves \eqref{eq:C-v96}. Diameters of two compact sets differ
by at most twice their Hausdorff distance. Multiplying
\eqref{eq:uniform-Hausdorff-v96} by $\sqrt t$ proves
\eqref{eq:uniform-law-v96}. Separation of the base clusters prevents
cross-cluster optimal matchings, uniformly on $H$. Continuity and
positivity follow from Lemma~\ref{lem:coercive-v96}.
\end{proof}

\begin{corollary}[Uniform local oracle risk]
\label{cor:uniform-risk-v96}
Fix $0<c\le1/(8\sqrt2)$ and $H\Subset\mathfrak I$. At the design of
each centre take independent categorical samples with $n_j/N\to w_j$
uniformly on $H$; integer rounding of $Nw_j$ is sufficient. For the
squared minimax risk $\mathcal R_N(\theta,c)$ on the specified ball
$h_w(F_T(\xi),P_\theta)\le cN^{-1/2}$, one has, uniformly on $H$,
\[
 {3c\over32}C(\theta)^2-O_H(N^{-1/4})
 \le\sqrt N\,\mathcal R_N(\theta,c)
 \le cC(\theta)^2+O_H(N^{-1/4}).
\]
The centre is part of this local decision problem; no unknown-centre
or global adaptation is asserted.
\end{corollary}
\begin{proof}
The constant-decision upper bound is $\omega_\theta(t_N)^2$.
For a diameter pair in the ball, the product Hellinger distance is at
most $\sqrt{8N}\,t_N\le1/4$, eventually uniformly on $H$.
Total variation is at most this distance. Testing the two targets
therefore gives the lower bound $3\omega_\theta(t_N)^2/32$.
Substitute $t_N=cN^{-1/2}$ into \eqref{eq:uniform-law-v96}; squaring
its uniform $O(t_N)$ remainder gives the displayed $O(N^{-1/4})$
error after multiplication by $\sqrt N$.
\end{proof}
